\scanpage{167}
\mspage{116}
% archive sheet number: 167
on trouve un foncteur

\[
\begin{aligned}
&\deleted{\unsure{\pi}}\\[-2pt]
i&\longmapsto\pi_0(X_i):\quad I\longrightarrow\mathrm{Ens}\\[-2pt]
&\hspace{3em}\tikz[baseline=-.4ex]{\node[rotate=90,inner sep=0pt] {$\simeq$};}\\[-2pt]
&\hspace{1em}\pi_0(X_{/i})
\end{aligned}
\]

Ceci posé, on a le résultat \unsure{subtil|crucial}

\uline{Proposition 12} : Soit $f:X\longrightarrow I$ $W$-propre, et supposons\textsuperscript{*} que $W$ satisfasse l'axiome de connexité $W(5)$. (sinon la $W$-propreté ne nous sera d'aucune utilité \ldots) On a alors une bijection canonique
% slanted
\margin{Il vaut mieux dire : Soit $f$ $W_0$-propre, où $W_0$ est le \unsure{localis\textsuperscript{t}} $0$-connexité \ldots{}}

\[
\pi_0(X)\xleftarrow{\sim}\underset{I}{\varinjlim}\,\pi_0(X_i),\tag*{\unsure{$(\alpha)$}}
\]
où la flèche est induite par les applications can. $\pi_0(X_i)\unsure{-|\to}\pi_0(X)$.

\uline{Dém} Passant à $f^\circ:X^\circ-I^\circ$, on est ramené à prouver le $=$ résultat pour $X$ \deleted{\unsure{lisse}}\uline{lisse} sur $I$, avec la seule différence que le foncteur

\[
i\longmapsto\pi_0(X_i)
\]
sur $I$ sera contravariant, i.e. un préfaisceau, ce qui offre l'avantage de l'intuition attachée au formalisme des préfaisceaux. Soit

\[
g:I\longrightarrow e,\qquad\text{d'où }\;gf=p:X-e,
\]
\deleted{\unsure{il}} est immédiat, par le calcul standard

\[
\deleted{\unsure{\pi}}\qquad p_!(F)=\pi_0(X_{/F})\qquad(F\in\mathrm{Ob}\,X^\wedge)
\]
\deleted{\unsure{que}} (on a identifié $e^\wedge$ à $(\mathrm{Ens})$), qu'on a

\[
\pi_0(X)=p_!(e_X)=g_!(f_!(e_X)).
\]

Or le préfaisceau

\[
i\longmapsto\pi_0(X_i)\quad(\simeq\pi_0(\unsure{{}_i\backslash X|i\backslash X}),\ \text{car }X\text{ \inserted{$W$-}lisse sur }I)
\]
